\documentclass{amsart}
% For dealing with references we use the comment environment
\usepackage{verbatim}
\newenvironment{reference}{\comment}{\endcomment}
%\newenvironment{reference}{}{}
\newenvironment{slogan}{\comment}{\endcomment}
\newenvironment{history}{\comment}{\endcomment}

% For commutative diagrams we use Xy-pic
\usepackage[all]{xy}

% We use 2cell for 2-commutative diagrams.
\xyoption{2cell}
\UseAllTwocells

% We use multicol for the list of chapters between chapters
\usepackage{multicol}

% This is generally recommended for better output
\usepackage{lmodern}
\usepackage[T1]{fontenc}

% For cross-file-references

% Package for hypertext links:
\usepackage{hyperref}

% For any local file, say "hello.tex" you want to link to please

% Theorem environments.
%
\theoremstyle{plain}
\newtheorem{theorem}[subsection]{Theorem}
\newtheorem{proposition}[subsection]{Proposition}
\newtheorem{lemma}[subsection]{Lemma}

\theoremstyle{definition}
\newtheorem{definition}[subsection]{Definition}
\newtheorem{example}[subsection]{Example}
\newtheorem{exercise}[subsection]{Exercise}
\newtheorem{situation}[subsection]{Situation}

\theoremstyle{remark}
\newtheorem{remark}[subsection]{Remark}
\newtheorem{remarks}[subsection]{Remarks}

\numberwithin{equation}{subsection}

% Macros
%
\def\lim{\mathop{\mathrm{lim}}\nolimits}
\def\colim{\mathop{\mathrm{colim}}\nolimits}
\def\Spec{\mathop{\mathrm{Spec}}}
\def\Hom{\mathop{\mathrm{Hom}}\nolimits}
\def\Ext{\mathop{\mathrm{Ext}}\nolimits}
\def\SheafHom{\mathop{\mathcal{H}\!\mathit{om}}\nolimits}
\def\SheafExt{\mathop{\mathcal{E}\!\mathit{xt}}\nolimits}
\def\Sch{\mathit{Sch}}
\def\Mor{\mathop{\mathrm{Mor}}\nolimits}
\def\Ob{\mathop{\mathrm{Ob}}\nolimits}
\def\Sh{\mathop{\mathit{Sh}}\nolimits}
\def\NL{\mathop{N\!L}\nolimits}
\def\CH{\mathop{\mathrm{CH}}\nolimits}
\def\proetale{{pro\text{-}\acute{e}tale}}
\def\etale{{\acute{e}tale}}
\def\QCoh{\mathit{QCoh}}
\def\Ker{\mathop{\mathrm{Ker}}}
\def\Im{\mathop{\mathrm{Im}}}
\def\Coker{\mathop{\mathrm{Coker}}}
\def\Coim{\mathop{\mathrm{Coim}}}

% Boxtimes
%
\DeclareMathSymbol{\boxtimes}{\mathbin}{AMSa}{"02}

%
% Macros for moduli stacks/spaces
%
\def\QCohstack{\mathcal{QC}\!\mathit{oh}}
\def\Cohstack{\mathcal{C}\!\mathit{oh}}
\def\Spacesstack{\mathcal{S}\!\mathit{paces}}
\def\Quotfunctor{\mathrm{Quot}}
\def\Hilbfunctor{\mathrm{Hilb}}
\def\Curvesstack{\mathcal{C}\!\mathit{urves}}
\def\Polarizedstack{\mathcal{P}\!\mathit{olarized}}
\def\Complexesstack{\mathcal{C}\!\mathit{omplexes}}
% \Pic is the operator that assigns to X its picard group, usage \Pic(X)
% \Picardstack_{X/B} denotes the Picard stack of X over B
% \Picardfunctor_{X/B} denotes the Picard functor of X over B
\def\Pic{\mathop{\mathrm{Pic}}\nolimits}
\def\Picardstack{\mathcal{P}\!\mathit{ic}}
\def\Picardfunctor{\mathrm{Pic}}
\def\Deformationcategory{\mathcal{D}\!\mathit{ef}}

\setlength{\emergencystretch}{3em}
\begin{document}
\title[Stacks Project, Tag 0AG0]{424 --- Unique factorization in regular local rings}
\maketitle
\noindent Copyright (C) 2005--2025 Johan de Jong. Stacks Project authors. Source: \href{https://stacks.math.columbia.edu/tag/0AG0}{Tag 0AG0}.

\noindent Editorial difficulty note (not author text). Difficulty H2: The proof inducts on dimension and makes a height-one prime invertible after inverting a regular parameter. Vanishing of the relevant Picard group supplies a generator there, which is descended using irreducible factorization and the behavior of prime elements under localization..

\noindent Editorial provenance note (not author text). History: excerpt from revision \texttt{a0444\allowbreak 6e57e\allowbreak c1fbc\allowbreak 252a8\allowbreak 71afc\allowbreak ec775\allowbreak 2fb28\allowbreak 07b14}, more-algebra.tex, lines 36800--36843. Reference presentation was adapted; original-author context and core support blocks were retained or added. The mathematical wording is unchanged. Licensed under GNU FDL 1.2 or later; no invariant sections or cover texts. See ../licenses/COPYING. Context blocks reproduce author definitions and supporting results; they are not additional counted problems.

\begin{lemma}
\label{lemma-regular-local-Pic-zero}
Let $R$ be a regular local ring. Let $f \in R$.
Then $\Pic(R_f) = 0$.
\end{lemma}

\begin{proof}
Let $L$ be an invertible $R_f$-module. In particular $L$ is
a finite $R_f$-module. There exists a finite $R$-module
$M$ such that $M_f \cong L$, see
Algebra, Lemma \href{https://stacks.math.columbia.edu/tag/05N5}{lemma construct fp module (Tag 05N5)}.
By Algebra, Proposition \href{https://stacks.math.columbia.edu/tag/00O7}{proposition regular finite gl dim (Tag 00O7)}
we see that $M$ has a finite free resolution $F_\bullet$ over $R$.
It follows that $L$ is quasi-isomorphic to a finite complex
of {\it free} $R_f$-modules. Hence by
Lemma \href{https://stacks.math.columbia.edu/tag/0AFY}{lemma perfect to K group (Tag 0AFY)} we see that
$[L_f] = n[R_f]$ in $K_0(R_f)$ for some $n \in \mathbf{Z}$.
Applying the map of Lemma \href{https://stacks.math.columbia.edu/tag/0AFX}{lemma det (Tag 0AFX)}
we see that $L$ is trivial.
\end{proof}

\begin{lemma}
\label{lemma-regular-local-UFD}
A regular local ring is a UFD.
\end{lemma}

\begin{proof}
Recall that a regular local ring is a domain, see
Algebra, Lemma \href{https://stacks.math.columbia.edu/tag/00NP}{lemma regular domain (Tag 00NP)}.
We will prove the unique factorization property
by induction on the dimension of the regular local ring $R$.
If $\dim(R) = 0$, then $R$ is a field and in particular a UFD.
Assume $\dim(R) > 0$. Let $x \in \mathfrak m$, $x \not \in \mathfrak m^2$.
Then $R/(x)$ is regular by Algebra, Lemma \href{https://stacks.math.columbia.edu/tag/00NQ}{lemma regular ring CM (Tag 00NQ)},
hence a domain by
Algebra, Lemma \href{https://stacks.math.columbia.edu/tag/00NP}{lemma regular domain (Tag 00NP)},
hence $x$ is a prime element.
Let $\mathfrak p \subset R$ be a height $1$ prime. We have
to show that $\mathfrak p$ is principal, see
Algebra, Lemma \href{https://stacks.math.columbia.edu/tag/0AFT}{lemma characterize UFD height 1 (Tag 0AFT)}.
We may assume $x \not \in \mathfrak p$, since if $x \in \mathfrak p$,
then $\mathfrak p = (x)$ and we are done.
For every nonmaximal prime $\mathfrak q \subset R$
the local ring $R_\mathfrak q$ is a regular local ring, see
Algebra, Lemma \href{https://stacks.math.columbia.edu/tag/0AFS}{lemma localization of regular local is regular (Tag 0AFS)}.
By induction we see that $\mathfrak pR_\mathfrak q$ is principal.
In particular, the $R_x$-module $\mathfrak p_x = \mathfrak pR_x \subset R_x$
is a finitely presented $R_x$-module whose localization at
any prime is free of rank $1$. 
By Algebra, Lemma \href{https://stacks.math.columbia.edu/tag/00NX}{lemma finite projective (Tag 00NX)}
we see that $\mathfrak p_x$ is an invertible $R_x$-module.
By Lemma \ref{lemma-regular-local-Pic-zero} we see that
$\mathfrak p_x = (y)$ for some $y \in R_x$.
We can write $y = x^e f$ for some $f \in \mathfrak p$ and $e \in \mathbf{Z}$.
Factor $f = a_1 \ldots a_r$ into irreducible elements of $R$
(Algebra, Lemma \href{https://stacks.math.columbia.edu/tag/034R}{lemma factorization exists (Tag 034R)}).
Since $\mathfrak p$ is prime, we see that $a_i \in \mathfrak p$
for some $i$. Since $\mathfrak p_x = (y)$ is prime and
$a_i | y$ in $R_x$, it follows that $\mathfrak p_x$ is generated by
$a_i$ in $R_x$, i.e., the image of $a_i$ in $R_x$ is prime.
As $x$ is a prime element, we find that $a_i$ is prime in $R$ by
Algebra, Lemma \href{https://stacks.math.columbia.edu/tag/0AFU}{lemma invert prime elements (Tag 0AFU)}.
Since $(a_i) \subset \mathfrak p$ and $\mathfrak p$ has height
$1$ we conclude that $(a_i) = \mathfrak p$ as desired.
\end{proof}
\end{document}
